\documentclass[11pt,a4paper]{article}
\usepackage[margin=25mm,headheight=14pt]{geometry}
\usepackage[T1]{fontenc}
\usepackage[utf8]{inputenc}
\usepackage{lmodern,amsmath,amssymb}
\usepackage{xcolor,graphicx,booktabs,longtable,array}
\usepackage{microtype,enumitem,fancyhdr,titlesec}
\usepackage{hyperref}
\definecolor{accent}{HTML}{163B55}
\hypersetup{colorlinks=true,linkcolor=accent,urlcolor=accent,
 pdftitle={A flat torus with logarithmic Sobolev constant below half its spectral gap},pdfauthor={Matthew J. Colbrook},
 pdfsubject={AIM 1019: mathematical result and proof}}
\urlstyle{same}
\Urlmuskip=0mu plus 1mu
\setlength{\parindent}{0pt}
\setlength{\parskip}{0.35em}
\setlength{\emergencystretch}{3em}
\setlist{itemsep=0.25em,topsep=0.3em}
\providecommand{\tightlist}{\setlength{\itemsep}{0pt}\setlength{\parskip}{0pt}}
\titleformat{\section}{\normalfont\large\bfseries\color{accent}}{\thesection}{0.7em}{}
\titleformat{\subsection}{\normalfont\normalsize\bfseries}{\thesubsection}{0.6em}{}
\pagestyle{fancy}
\fancyhf{}
\fancyhead[L]{\small\sffamily AIM 1019}
\fancyhead[R]{\small\sffamily Matthew J. Colbrook}
\fancyfoot[C]{\thepage}
\renewcommand{\headrulewidth}{0.3pt}
\setcounter{secnumdepth}{0}
\allowdisplaybreaks[2]
\title{\textbf{A flat torus with logarithmic Sobolev constant below half its spectral gap}}
\author{Matthew J. Colbrook\\[3pt]
{\small Department of Applied Mathematics and Theoretical Physics}\\
{\small University of Cambridge, United Kingdom}\\
{\small\href{mailto:mjc249@maths.cam.ac.uk}{\texttt{mjc249@maths.cam.ac.uk}}}}
\date{3 November 2026}
\begin{document}
\maketitle
\begin{abstract}
We exhibit a flat torus with spectral gap one and logarithmic Sobolev constant at most $1296/2669<1/2$. An explicit positive trial function and a rational entropy estimate prove the strict inequality. The example extends to an open family of metrics and to higher dimensions; the general sharp-constant question in AIM 1019 remains open.
\end{abstract}
\noindent\textbf{Repository target.}
This manuscript concerns AIM problem 1019 at revision
\texttt{3fa50884d129}. Its scope is stated below.
\medskip

Let \(T\) be a flat torus with normalized volume measure \(\mu\). Write

\[
\mathcal E(f,f)=\int_T|\nabla f|^2\,d\mu,
\]

\[
\operatorname{Ent}_\mu(f^2)
=\int_T f^2\log f^2\,d\mu
-\left(\int_T f^2\,d\mu\right)
 \log\left(\int_T f^2\,d\mu\right),
\]

and define

\[
\alpha(T)=\inf_{\substack{f\in H^1(T;\mathbb R)\\
\operatorname{Ent}_\mu(f^2)>0}}
\frac{\mathcal E(f,f)}{\operatorname{Ent}_\mu(f^2)}.
\]

Let \(\lambda(T)\) denote the first positive Laplace eigenvalue, or
spectral gap.

The universal equality proposed in Problem 1019 is false. We give an
explicit flat two-dimensional torus with spectral gap one and prove

\[
\boxed{\alpha(T)\leq\frac{1296}{2669}<\frac12=\frac{\lambda(T)}2.}
\]

The argument uses the classical third-moment obstruction to equality,
together with a rational estimate that makes the strict inequality
quantitative. The obstruction is stated, for example, in
Chen--Liu--Saloff-Coste, Theorem 1.1. The construction also gives an
open family of counterexamples and extends to every dimension at least
two.

This does not determine the exact sharp logarithmic Sobolev constant,
even for the displayed torus. It therefore resolves the
universal-equality part of the catalogue entry, while leaving its
request to determine the sharp constant for general lattices unanswered.

\hypertarget{1019-the-torus-and-a-resonant-first-eigenfunction}{%
\section{1. The torus and a resonant first
eigenfunction}\label{1019-the-torus-and-a-resonant-first-eigenfunction}}

Use angular coordinates \((x,y)\in\mathbb R^2/(2\pi\mathbb Z)^2\) and
the constant metric whose inverse matrix is

\[
G=\begin{pmatrix}1&-1/2\\-1/2&1\end{pmatrix}.
\]

The normalized volume measure is \(d\mu=dx\,dy/(2\pi)^2\). The positive
Laplacian acts on a Fourier mode by

\[
-\Delta e^{i(mx+ny)}=q(m,n)e^{i(mx+ny)},\qquad
q(m,n)=m^2-mn+n^2.
\]

For a nonzero integer pair \((m,n)\), this positive integer is at least
one. Equality holds at \((1,0)\), \((0,1)\), \((1,1)\), and their
negatives. Thus the spectral gap is

\[
\lambda(T)=1.
\]

Equivalently, the dual Euclidean frequency lattice is generated by

\[
a=(1,0),\qquad b=(-1/2,\sqrt3/2),
\]

so \(a,b,a+b\) all have length one. A direct Euclidean lattice defining
the same torus is generated by

\[
(2\pi,2\pi/\sqrt3),\qquad (0,4\pi/\sqrt3).
\]

Set

\[
g(x,y)=\cos x+\cos y+\cos(x+y).
\]

Then \(-\Delta g=g\). Fourier orthogonality gives

\[
\int_T g\,d\mu=0,\qquad
\int_T g^2\,d\mu=\frac32,\qquad
\int_T g^3\,d\mu=\frac32,\qquad
\int_T g^4\,d\mu=\frac{45}{8}.
\tag{1}
\]

The cubic moment comes from the term \(6\cos x\cos y\cos(x+y)\), whose
integral is \(6/4\). For the fourth moment, there are \(90\) ordered
four-step closed walks with steps in \(\{\pm a,\pm b,\pm(a+b)\}\);
dividing by \(2^4\) gives \(45/8\). The same value follows by expanding
the fourth power and using Fourier orthogonality.

We will also use the range bound

\[
-\frac32\leq g\leq3.
\tag{2}
\]

Indeed, writing \(s=(x+y)/2\) and \(t=(x-y)/2\) gives

\[
g=2\cos s\cos t+2\cos^2s-1.
\]

For fixed \(r=|\cos s|\), its minimum is \(2r^2-2r-1\geq-3/2\). The
upper bound is immediate.

\hypertarget{1019-the-third-moment-obstruction}{%
\section{2. The third-moment
obstruction}\label{1019-the-third-moment-obstruction}}

Take \(f_\varepsilon=1+\varepsilon g\). Since \(g\) is a first
eigenfunction,

\[
\mathcal E(f_\varepsilon,f_\varepsilon)=\frac32\varepsilon^2.
\]

Taylor expansion, uniform because \(g\) is bounded, gives

\[
\begin{aligned}
\operatorname{Ent}_\mu(f_\varepsilon^2)
&=2\varepsilon^2\int_T g^2\,d\mu
 +\frac23\varepsilon^3\int_T g^3\,d\mu
 +O(\varepsilon^4)\\
&=3\varepsilon^2+\varepsilon^3+O(\varepsilon^4).
\end{aligned}
\]

Consequently the energy-to-entropy ratio is strictly less than \(1/2\)
for all sufficiently small positive \(\varepsilon\). This already proves

\[
\alpha(T)<\lambda(T)/2.
\]

\hypertarget{1019-a-rational-certificate-of-strict-inequality}{%
\section{3. A rational certificate of strict
inequality}\label{1019-a-rational-certificate-of-strict-inequality}}

Choose \(\varepsilon=1/6\) and \(f=1+g/6\). By (2), \(f\geq3/4\). Define

\[
h(z)=(1+z)^2\log((1+z)^2).
\]

Its Taylor polynomial through degree three is \(2z+3z^2+(2/3)z^3\), and

\[
h^{(4)}(z)=-\frac4{(1+z)^2}.
\]

Between zero and any \(z=g/6\), one has \(1+z\geq3/4\). Taylor's theorem
therefore yields

\[
h(z)\geq2z+3z^2+\frac23z^3-\frac8{27}z^4.
\tag{3}
\]

For \(s\geq0\), one also has

\[
(1+s)\log(1+s)\leq s+\frac{s^2}{2}.
\tag{4}
\]

For example, the derivative of the difference between the right and left
sides is \(s-\log(1+s)\geq0\), and the difference vanishes at zero.

Here

\[
\int_T f^2\,d\mu=1+\frac{3/2}{36}=\frac{25}{24}.
\]

Combining (1), (3), and (4) gives

\[
\begin{aligned}
\operatorname{Ent}_\mu(f^2)
&\geq3\varepsilon^2+\varepsilon^3
-\varepsilon^4\left(
\frac{45/8}{6(3/4)^2}+\frac{(3/2)^2}{2}\right)\\
&=\frac1{12}+\frac1{216}-\frac{67}{31104}
=\frac{2669}{31104}.
\end{aligned}
\]

Since \(\mathcal E(f,f)=1/24\), this proves

\[
\boxed{\alpha(T)\leq\frac{1/24}{2669/31104}
=\frac{1296}{2669}<\frac12.}
\tag{5}
\]

The bound in (5) is a strict upper certificate; it is not asserted to be
the sharp constant.

\hypertarget{1019-an-open-family-of-counterexamples}{%
\section{4. An open family of
counterexamples}\label{1019-an-open-family-of-counterexamples}}

Replace \(G\) by

\[
G_c=\begin{pmatrix}1&-c\\-c&1\end{pmatrix},\qquad0\leq c<1,
\]

and denote the resulting torus by \(T_c\). The normalized measure, the
function \(g\), and the entropy of \(f=1+g/6\) are unchanged. Its energy
is

\[
\mathcal E_c(f,f)=\frac{2-c}{36}.
\]

The spectral gap is

\[
\lambda(T_c)=\min\{1,2(1-c)\}.
\tag{6}
\]

To verify (6), replacing integers \(m,n\) by their absolute values
cannot increase \(m^2-2cmn+n^2\). If one coordinate is zero, the value
is at least one. Otherwise take \(m=n+r\) with \(n\geq1\) and
\(r\geq0\). Then

\[
m^2-2cmn+n^2=2(1-c)n(n+r)+r^2,
\]

which is at least \(2(1-c)\) when \(r=0\) and at least one when
\(r\geq1\). Both candidate minima are attained.

The same trial function therefore gives

\[
\alpha(T_c)\leq\frac{864(2-c)}{2669}.
\]

Comparison with (6) shows that this is strictly smaller than
\(\lambda(T_c)/2\) throughout the explicit interval

\[
\boxed{\frac{787}{1728}<c<\frac{941}{1805}.}
\]

This interval contains \(c=1/2\). More generally, sufficiently small
changes in all metric coefficients preserve the strict inequality for
this fixed witness, by continuity of its energy and of the shortest
dual-lattice length. Thus failure of equality is not confined to a
single lattice with extra triangular-lattice multiplicity.

\hypertarget{1019-higher-dimensions-and-the-remaining-sharp-constant-question}{%
\section{5. Higher dimensions and the remaining sharp-constant
question}\label{1019-higher-dimensions-and-the-remaining-sharp-constant-question}}

For \(d>2\), take the product of \(T\) with \(d-2\) standard circles of
circumference \(2\pi\). The product spectral gap is still one. The same
function \(f\), independent of the extra coordinates, has the same
energy and entropy with respect to normalized volume. Hence the
universal equality fails in every dimension \(d\geq2\).

The exact value of \(\alpha\) for the triangular torus, and a
determination of \(\alpha\) from an arbitrary lattice, are not supplied
by this argument.

\hypertarget{1019-references-and-scope}{%
\section{References and scope}\label{1019-references-and-scope}}

\begin{enumerate}
\def\labelenumi{\arabic{enumi}.}
\tightlist
\item
  G.-Y. Chen, W.-W. Liu, and L. Saloff-Coste, \emph{The logarithmic
  Sobolev constant of some finite Markov chains}, Annales de la Faculté
  des Sciences de Toulouse \textbf{17} (2008), 239--290. Theorem 1.1
  gives the classical third-moment obstruction; Question Q1 on printed
  p.~242 asks about the flat-torus equality.
  \url{https://www.numdam.org/item/10.5802/afst.1183.pdf}.
\item
  L. Saloff-Coste, \emph{Analysis on compact Lie groups of large
  dimension and on connected compact groups}, Colloquium Mathematicum
  \textbf{118} (2010), 183--199, Section 3.4. Its Problem 11 is a
  comparison question, not the exact equality, and is not used as a
  source for that equality.
  \url{https://pi.math.cornell.edu/~lsc/papers/hulancoll.pdf}.
\end{enumerate}
\end{document}
